% Fichier généré automatiquement par tools/generate_corrections.py.
% Source : PERF/PERF-02-Marges/63_BancHydraulique/63_BancHydraulique.tex
% Statut : complet (7/7 question(s) renseignée(s), 0 reprise(s)).
\begin{corrigebox}[Corrigé extrait de la banque]
\CorrigeQuestion{1}
On a $C(p)=\dfrac{K_i}{p}+K_p = \dfrac{K_i+p K_p}{p} = K_i \dfrac{1+p \dfrac{K_p}{K_i}}{p}$.
\CorrigeQuestion{2}
\begin{marginfigure}
\includegraphics[width=\linewidth]{63_02_cor}
\end{marginfigure}

\begin{marginfigure}
\includegraphics[width=\linewidth]{63_03_cor}
\end{marginfigure}
\CorrigeQuestion{3}
Ce correcteur augmente la classe de la FTBO donc augmente la précision. 
Cependant, il réduit la phase. Il faut donc veiller à ce que la pulsation de cassure soit réglée de telle sorte que le système ne soit pas déstabilisé.
\CorrigeQuestion{4}
D'après la remarque, on a $t_e \omega_{\SI{0}{dB}}=3$ soit $\omega_{\SI{0}{dB}} =3/t_e = \SI{0,075}{rad.s^{-1}}$.
\CorrigeQuestion{5}
Calculons la fonction de transfert en boucle ouverte non corrigée : 
$\indice{F}{BO}=\dfrac{\indice{K}{pom}}{1+T_2 p}\dfrac{\indice{K}{m}}{1+T_1 p} \indice{K}{cap}$.

Le correcteur doit être réglé pour que $\omega_{\SI{0}{dB}} =\SI{0,075}{rad.s^{-1}}$. 

Calculons la marge de phase. 
$\arg\left(\indice{F}{BO}\right)=-\arg\left(1+T_1p\right)-\arg\left(1+T_2p\right) = -\arctan T_1 \omega -\arctan T_2 \omega $
+
On a donc $\arg\left(\indice{F}{BO}(0,075)\right) = -\arctan (10 \times  0,075) -\arctan (5 \times  0,075)=-57\degres$ soit une marge de phase de $-123\degres$.

Pour atteindre une marge de phase de 60\degres, on peut donc baisser la phase de 63\degres.

Calculons $\arg\left(C(j\omega)\right)=-90+\arctan \left(\dfrac{K_p}{K_i}\omega\right)$.

On cherche donc $\dfrac{K_i}{K_p}$ tel que $\arg\left(C(0,075)\right)=-63$
Soit 
$-90+\arctan \left(\dfrac{K_p}{K_i}0,075\right) =-63$
$\Leftrightarrow \arctan \left(\dfrac{K_p}{K_i}0,075\right) =27$
$\Rightarrow  \dfrac{K_p}{K_i}0,075 =0,51$
$\Leftrightarrow  \dfrac{K_p}{K_i}=6,79$.
\CorrigeQuestion{6}
Il faut chercher $K_i$ et $K_p$ pour respecter $\omega_{\SI{0}{dB}}$. Recherchons le gain de la boucle ouverte non corrigée pour $\omega_{\SI{0}{dB}}$.

$\indice{G}{dB}\left(\indice{F}{BO}\right)=20\log\left(\indice{K}{pom}\indice{K}{m}\indice{K}{cap}\right)
-20\log\left(\sqrt{1^2+T_1^2\omega^2}\right)-20\log\left(\sqrt{1^2+T_2^2\omega^2}\right)$

On a alors $\indice{G}{dB}\left(\indice{F}{BO}\right)(0,075) = -0,004-1,94-0,57=\SI{2,52}{dB}$.

Il faut donc baisser le gain de \SI{2,52}{dB}
$\indice{G}{dB}\left(C(p)\right)$  
$= 20\log K_i -20\log \omega +20\log\left(\sqrt{1+\left(\dfrac{K_p}{K_i}\right)^2\omega^2}\right) $.

On a alors $\indice{G}{dB}\left(C(0,075)\right) = 20\log K_i +22,5+1 = -2,52$ soit $K_i = 10^{-\dfrac{2,52+1+22,5}{20}  } =0,05$.

Par suite, $K_p = 6,79 \times 0,05 = 0,34$.

(\textbf{A vérifier}).
\CorrigeQuestion{7}
\begin{marginfigure}
\includegraphics[width=\linewidth]{63_04_cor}
\end{marginfigure}
\begin{itemize}
\item Stabilité : 
\begin{itemize}
\item Marge de phase mesurée :  60\degres \textbf{cdc ok}.
\item Marge de gain mesurée :   infini \textbf{cdc ok}.
\end{itemize}
\item Rapidité : $ t_e = \SI{32}{s} < \SI{40}{s}$ \textbf{cdc ok}.
\item Précision :  écart statique nul \textbf{cdc ok}.
\item Amortissement :  nul \textbf{cdc ok}.
\end{itemize}

\begin{enumerate}
  \item $C(p)= K_i \dfrac{1+p \dfrac{K_p}{K_i}}{p}$.
  \item .
  \item .
  \item $T = 6,79$.
  \item $K_i = 0,05$ et $K_p=0,34$ (à vérifier).
\end{enumerate}
\end{corrigebox}
